% ============================================================
%  Rapport de projet d'ingénierie -- ShootingGameVRAR
% ============================================================
\documentclass[12pt, a4paper]{article}

% ---------- Encodage & langue ----------
\usepackage[utf8]{inputenc}
\usepackage[T1]{fontenc}
\usepackage[french]{babel}

% ---------- Mise en page ----------
\usepackage[
  top=2.5cm, bottom=2.5cm,
  left=2.5cm, right=2.5cm,
  headheight=15pt
]{geometry}
\usepackage{fancyhdr}       % En-têtes et pieds de page personnalisés
\usepackage{emptypage}      % Pages vierges vraiment vierges
\usepackage{setspace}       % Interligne
\setstretch{1.15}

% ---------- Typographie ----------
\usepackage{lmodern}        % Police Latin Modern (qualité vectorielle)
\usepackage{microtype}      % Optimisation microtypographique
\usepackage{csquotes}       % Guillemets typographiques français

% ---------- Couleurs & graphismes ----------
\usepackage[dvipsnames, table]{xcolor}
\definecolor{primary}{RGB}{0, 82, 155}      % Bleu ingénierie
\definecolor{secondary}{RGB}{230, 240, 255} % Fond léger
\definecolor{codebg}{RGB}{245, 245, 245}    % Fond blocs de code
\definecolor{accent}{RGB}{220, 50, 50}      % Rouge accent

\usepackage{graphicx}       % Insertion d'images
\usepackage{tikz}           % Dessins vectoriels
\usepackage{float}          % Contrôle fin du placement des flottants

% ---------- Mathématiques ----------
\usepackage{amsmath, amssymb, amsthm}
\usepackage{mathtools}

% ---------- Tableaux ----------
\usepackage{booktabs}       % Filets professionnels (\toprule, \midrule...)
\usepackage{tabularx}       % Tableaux a largeur fixe
\usepackage{multirow}       % Fusion de cellules

% ---------- Listes ----------
\usepackage{enumitem}
\setlist[itemize]{label=\textcolor{primary}{\textbullet}, itemsep=4pt}
\setlist[enumerate]{itemsep=4pt}

% ---------- Code source ----------
\usepackage{listings}
\lstset{
  backgroundcolor=\color{codebg},
  basicstyle=\ttfamily\small,
  keywordstyle=\color{primary}\bfseries,
  commentstyle=\color{gray}\itshape,
  stringstyle=\color{accent},
  numbers=left, numberstyle=\tiny\color{gray},
  numbersep=8pt,
  frame=single, framerule=0.4pt, rulecolor=\color{primary!40},
  breaklines=true, captionpos=b,
  tabsize=4, showstringspaces=false
}

% ---------- References & liens ----------
\usepackage[
  colorlinks=true,
  linkcolor=primary,
  urlcolor=primary,
  citecolor=primary,
  pdftitle={Rapport -- ShootingGameVRAR},
  pdfauthor={Autuly, Pelletier, Dje, Kouassi},
  pdfsubject={Projet d'ingenierie VR/AR},
  pdfkeywords={VR, AR, Unity, Meta Quest 2, Android, jeu de tir}
]{hyperref}

% ---------- Boites encadrees ----------
\usepackage[most]{tcolorbox}
\tcbset{
  colback=secondary,
  colframe=primary,
  fonttitle=\bfseries,
  boxrule=0.6pt,
  arc=4pt
}

% ---------- En-tetes / pieds de page ----------
\pagestyle{fancy}
\fancyhf{}
\fancyhead[L]{\small\textcolor{primary}{\textbf{ShootingGameVRAR}}}
\fancyhead[R]{\small\textcolor{gray}{Rapport de projet}}
\fancyfoot[C]{\textcolor{gray}{\thepage}}
\renewcommand{\headrulewidth}{0.5pt}
\renewcommand{\headrule}{%
  \hbox to\headwidth{\color{primary}\leaders\hrule height \headrulewidth\hfill}}

% ---------- Style des sections ----------
\usepackage{titlesec}
\titleformat{\section}
  {\Large\bfseries\color{primary}}
  {\thesection.}{0.6em}{}
  [\vspace{-0.5em}\textcolor{primary!40}{\rule{\linewidth}{0.5pt}}\vspace{0.2em}]
\titleformat{\subsection}
  {\large\bfseries\color{primary!80}}
  {\thesubsection.}{0.5em}{}
\titleformat{\subsubsection}
  {\normalsize\bfseries\color{primary!60}}
  {\thesubsubsection.}{0.5em}{}

% =============================================================
\begin{document}

% ---- Page de titre ----
\begin{titlepage}
  \pagecolor{primary}
  \color{white}
  \vspace*{2cm}
  \begin{center}
    {\large\scshape Rapport de Projet}\par
    \vspace{0.5cm}
    {\color{white!60}\rule{0.5\linewidth}{0.6pt}}\par
    \vspace{1.2cm}
    {\fontsize{30}{36}\selectfont\bfseries ShootingGameVRAR}\par
    \vspace{0.5cm}
    {\large Jeu de tir coopératif asymétrique VR\,/\,AR}\par
    \vspace{0.8cm}
    {\color{white!60}\rule{0.5\linewidth}{0.6pt}}\par
    \vspace{1.5cm}
    \href{https://github.com/TeodoreAutuly/ShootingGameVRAR/tree/main}{%
      \textcolor{white!80}{\small\texttt{github.com/TeodoreAutuly/ShootingGameVRAR}}}
  \end{center}
  \vfill
  \begin{minipage}{0.55\linewidth}
    \textbf{Auteurs~:}\\[6pt]
    Téodore \textsc{Autuly}\\
    Lise \textsc{Pelletier}\\
    Jean-Louis \textsc{Dje}\\
    Emmanuel \textsc{Kouassi}
  \end{minipage}
  \hfill
  \begin{minipage}{0.35\linewidth}\raggedleft
    \textbf{Date~:}\\[6pt]
    Mars 2026
  \end{minipage}
  \vspace{1.5cm}
\end{titlepage}
\nopagecolor


\tableofcontents
\newpage

% =============================================================
\section{Initialisation du projet}

\subsection{Concept de l'application}

Ce projet propose une expérience de tir coopérative reposant sur un \textbf{gameplay asymétrique} entre deux joueurs. Le premier, équipé d'un casque
\textbf{Meta Quest~2}, est immergé dans un environnement forestier o\`u il doit tirer
sur des cibles. Le second, utilisant un \textbf{smartphone Android}, agit comme un
drone éclaireur capable de se déplacer librement dans l'espace de jeu.

Sa mission est de détecter et d'activer des cibles invisibles au joueur VR, afin de
permettre au tireur de les voir et de les éliminer. L'expérience repose sur un défi
de communication et de synchronisation en temps réel entre les deux protagonistes.

\subsection{Interactions principales entre les joueurs}

On distingue deux grandes catégories d'interactions au sein de l'application :

\begin{itemize}
  \item \textbf{Interaction avec les cibles~:} le joueur AR \textbf{détecte} et \textbf{active} les
        cibles invisibles pour le joueur VR~; ce dernier doit ensuite \textbf{tirer} sur les
        cibles rendues visibles.

  \item \textbf{Interaction entre les joueurs~:} le joueur AR est représenté par un
        drone de combat dans l'environnement VR. Il \textbf{émet des signaux sonores}
        lors de l'activation de cibles, que le joueur VR perçoit. Réciproquement, le
        joueur VR \textbf{génère des signaux sonores en tirant}, per\c{c}us par le joueur AR. Une
        \textbf{communication libre} entre les deux joueurs est également possible.
\end{itemize}

\subsection{Architecture envisagée}

Nous envisageons les choix d'architectuire suivants pour le développement de l'application :
\begin{itemize}
  \item \textbf{Réseau et communication multijoueur :} nous utiliserons \textbf{Unity Netcode for GameObjects} pour gérer la synchronisation en temps réel entre les deux joueurs, avec un modèle distribué.
  \item \textbf{Architecture :} nous adopterons une architecture \textbf{Modèle, Vue, Contrôleur (MVC)} pour séparer les différentes responsabilités au sein de l'application, facilitant ainsi la maintenance et l'évolution du code. Des \textbf{Services} et des \textbf{Configurations} permettront de centraliser les données partagées et les paramètres de jeu, tandis que l'application sera lancée à partir d'un \textbf{Bootstrap} qui initialisera les composants nécessaires.
\end{itemize}


\section{Synchronisation réseau et environnement}

\subsection{Évolution de l'application}

Le prototype a récemment évolué pour intégrer des mécaniques asymétriques plus poussées et une véritable séparation des rôles. Le projet s'articule aujourd'hui autour des espaces suivants :

\begin{itemize}
  \item \textbf{La Scène de Lobby} : Elle remplace l'ancienne interface de connexion basique. Le rôle de l'appareil est détecté automatiquement. Le joueur VR (Hôte) héberge la partie via Unity Relay et génère un "Join Code" affiché au sein d'un environnement 3D. Le joueur AR (Client), quant à lui, dispose d'une interface en 2D sur l'écran de son smartphone pour saisir ce code et rejoindre la session.
  
  \item \textbf{La Scène de Jeu} : Une fois les deux joueurs connectés, le serveur charge automatiquement cette scène. Le joueur VR évolue dans un environnement, tandis que le joueur AR visionne le monde réel et voit les éléments du jeu en surcouche. Ce dernier est désormais représenté au niveau réseau par le modèle 3D d'un drone. Les contrôles de ce drone sont assurés par un joystick tactile virtuel pour les déplacements, et la rotation spatiale est synchronisée physiquement de façon dynamique selon l'orientation de la caméra du smartphone grâce au script (\textit{ARDroneController}).
\end{itemize}

\subsection{Calibration spatiale par ancrage QR Code}

Une calibration entre l'environnement en VR et celui en AR est réalisée, nous avons implémenté un système d'ancrage s'appuyant sur \textbf{AR Foundation} :
\begin{itemize}
  \item \textbf{Côté VR (Serveur) :} un repère virtuel (\textit{marker}) désigne explicitement le point spatial fixe de calibration (l'origine de la carte).
  \item \textbf{Côté AR (Client) :} le joueur dispose initialement de son propre référentiel. Lorsqu'il scanne un QR code physique imprimé et déposé au sol (faisant office d'origine dans le monde réel), le système \textit{ARTrackedImageManager} calcule l'écart de distance (Offset de position) et d'orientation entre ce repère physique et le repère virtuel désigné sur le serveur.
  \item \textbf{Alignement :} ces décalages sont ensuite appliqués à l'espace global (\textit{AR Session Origin}) du joueur AR. En conséquence, l'environnement physique et le champ virtuel fusionnent numériquement : le drone AR apparaît toujours à la position exacte et relative dans le casque VR lorsque l'utilisateur Android se déplace physiquement de l'autre côté du QR Code.
\end{itemize}

\subsection{Architecture de synchronisation}

\begin{itemize}
  \item \textbf{Architecture Host-Client :} Le casque VR agit en tant qu'Hôte (il est à la fois le Serveur qui valide les actions et un Client qui joue), tandis que le téléphone AR se connecte en tant que Client.

  \item \textbf{Services Cloud Unity :} La connexion passe désormais par un relai internet via Unity Relay Services. Cette approche gère l'allocation des réseaux virtuels et des paquets de manière invisible et automatique une fois authentifié.

  \item \textbf{Netcode for GameObjects (NGO) :} Chaque joueur (Prefab VR et Prefab Drone AR) ainsi que les objets interactifs (cibles environnementales) possèdent un composant \textit{NetworkObject}, leur attribuant un identifiant unique.

  \item \textbf{Synchronisation de l'état :} Le composant \textit{NetworkTransform} réplique en temps réel la position et la rotation des joueurs. En outre, la manipulation et la distribution des états coopératifs (comme l'activation ou la destruction des cibles) s'appuie désormais sur des \textit{NetworkVariables} et requêtes \textit{ServerRpc}/\textit{ClientRpc}, ce qui garantit les échanges sur le réseau.
\end{itemize}

\section{Limites actuelles et défis techniques}

\subsection{Environnement 3D en cours de développement}

Le codage s'étant principalement concentré sur la validation asymétrique (le réseau, le système logique Lobby/Partie, et la configuration d'AR Foundation), l'intégration de l'environnement 3D complet demeure inachevée.

Le modèle détaillé d'une forêt, le design définitif des "cibles à détruire", ainsi que le maillage graphique du drone sont encore au stade préliminaire. Actuellement, l'application fonctionne avec du contenu volontairement basique. Néanmoins, nous avons trouvé des pistes de préfabriqués particulièrement adaptés à notre projet sur le Unity AssetStore.

\subsection{Difficultés de synchronisation}

En amont de la transition vers Unity Relay, l'un des premiers défis techniques fut la mise en place d'une connectivité réseau simple entre les deux terminaux. Au début du projet, l'interconnexion entre le casque VR et le smartphone AR demandait une initialisation manuelle à base d'adresses IP à entrer manuellement sur Unity Inspector à chaque fois.

Cela était une nuisance particulièrement significative pour l'ergonomie de l'application et nous avons décidé d'implémenter définitivement \textbf{Unity Relay} sans adresse IP manuelle, au profit d'un gameplay fluide.

\section{Interactions collaboratives}

\subsection{Description des interactions}

La nouvelle scène de \textit{Lobby} implémente désormais une logique de \textbf{menu local} propre à chaque appareil, en remplacement du lobby partagé utilisé au début du projet. Avant la connexion réseau, chaque joueur se trouve dans son propre contexte d'affichage et ne voit que les éléments nécessaires à son rôle.

\begin{itemize}
  \item \textbf{Côté VR (Meta Quest 2 -- Hôte)} : un menu 3D minimal est affiché dans l'espace du casque. Le joueur lance l'hébergement de session, ce qui initialise Unity Services, crée l'allocation Relay puis démarre l'instance \textit{Host} via Netcode for GameObjects.
  \item \textbf{Côté AR (Android -- Client)} : une interface 2D tactile présente un champ de saisie du \textit{Join Code} et un bouton de connexion. Après validation du code, le client rejoint l'allocation Relay et démarre sa session réseau.
\end{itemize}

Ce découplage clarifie le parcours utilisateur : chaque joueur interagit d'abord avec une interface adaptée à son matériel, puis la transition vers la scène commune est pilotée par le serveur lorsque la connexion est validée.

\subsection{Problèmes rencontrés et ajustements}

Le passage de l'ancien lobby vers ce nouveau menu a nécessité plusieurs ajustements techniques :

\begin{itemize}
  \item \textbf{Suppression de la dépendance au lobby partagé initial :} l'ancienne approche imposait des éléments visuels communs avant la connexion, ce qui créait des incohérences entre VR et mobile.
  \item \textbf{Réorganisation de l'initialisation réseau :} l'ordre d'appel (authentification, allocation Relay, démarrage NGO) a été fiabilisé pour éviter les échecs de connexion intermittents.
  \item \textbf{Sécurisation de la transition de scène :} le chargement de la scène de jeu est déclenché uniquement par l'hôte lorsque les deux rôles sont effectivement présents, afin d'éviter les désynchronisations.
  \item \textbf{Adaptation UX par plateforme :} une même interface n'était pas ergonomique pour les deux dispositifs ; le menu a donc été spécialisé (VR : menu spatial, AR : panneau tactile).
\end{itemize}

Ces ajustements ont permis de rendre l'entrée en partie plus lisible, plus stable et plus rapide à exécuter lors des tests.

\subsection{Environnement}

L'environnement de jeu reste centré sur une zone forestière servant de terrain commun aux deux joueurs après connexion. Le joueur VR y évolue en immersion complète, tandis que le joueur AR y projette les éléments virtuels via la calibration spatiale sur QR Code.

\begin{figure}[H]
  \centering
  \fbox{\rule{0pt}{6cm}\rule{0.92\linewidth}{0pt}}
  \caption{Espace image -- Vue d'ensemble de l'environnement de jeu (forêt, zones de cibles, repères de navigation).}
  \label{fig:env-overview}
\end{figure}

\begin{figure}[H]
  \centering
  \fbox{\rule{0pt}{6cm}\rule{0.92\linewidth}{0pt}}
  \caption{Espace image -- Représentation du drone AR dans l'environnement vu côté VR.}
  \label{fig:env-drone-vr}
\end{figure}

\subsection{Scénario de l'application}

Le déroulement collaboratif suit la séquence suivante :

\begin{enumerate}
  \item Le joueur VR ouvre l'application et héberge une partie depuis son menu local.
  \item Un \textit{Join Code} est généré et communiqué au joueur AR.
  \item Le joueur AR saisit ce code dans son interface mobile puis rejoint la session.
  \item Une fois les deux joueurs connectés, l'hôte déclenche le passage automatique à la scène commune.
  \item En partie, le joueur AR localise et active les cibles ; le joueur VR reçoit l'information visuelle/sonore et tire pour les détruire.
\end{enumerate}

\textbf{Détail de configuration (affichage des cibles).} Côté serveur, le script \textit{TargetManager} instancie un \textit{targetPrefab} sur des \textit{spawnPoints} prédéfinis, puis synchronise chaque cible avec un \textit{NetworkObject}. Chaque cible démarre non activée (variable réseau \textit{isActivated = false}) : elle est donc non tirable et visuellement neutre. Quand le joueur AR l'active via \textit{ServerRpc}, l'état est propagé à tous les clients par \textit{NetworkVariable}, le rendu passe en rouge et le collider est activé ; elle devient alors visible/tirable pour le joueur VR. Après impact validé côté serveur, la destruction est diffusée via \textit{ClientRpc} puis la cible est retirée du réseau.

Ce scénario met en valeur la complémentarité des rôles : l'AR assiste et révèle, la VR engage et neutralise.

\subsection{Interfaces UX}

L'UX a été pensée selon deux axes : \textbf{lisibilité immédiate} et \textbf{cohérence au rôle}. Le menu VR privilégie des actions peu nombreuses et visibles en 3D, tandis que le menu AR concentre la saisie du code et les actions de connexion sur des composants tactiles simples.

\begin{figure}[H]
  \centering
  \fbox{\rule{0pt}{6cm}\rule{0.92\linewidth}{0pt}}
  \caption{Espace image -- Interface du menu VR (création de session et affichage du code).}
  \label{fig:ux-vr-menu}
\end{figure}

\begin{figure}[H]
  \centering
  \fbox{\rule{0pt}{6cm}\rule{0.92\linewidth}{0pt}}
  \caption{Espace image -- Interface du menu AR (saisie du Join Code et bouton de connexion).}
  \label{fig:ux-ar-menu}
\end{figure}

\begin{figure}[H]
  \centering
  \fbox{\rule{0pt}{6cm}\rule{0.92\linewidth}{0pt}}
  \caption{Espace image -- Interface en jeu (repères d'activation des cibles et feedback collaboratif).}
  \label{fig:ux-ingame}
\end{figure}

\end{document}